%! Two Categorical Variables — Unit 2, Lesson 2.1 — Warm-Up (Answer Key)
% Mirrors the blank exactly apart from the package swap, the pageheader, and the answers.
% NO teachernote here — teacher prose lives in the lesson plan.
% ONE PAGE, blank and key. Three items at 12pt.
% GRADUAL RELEASE: the warm-up activates prior knowledge before the guided notes.
%   Item 1 spirals to Lesson 1.2 (relative frequency = count / total) — the notes' joint and
%          marginal relative frequencies are this move with a bigger table.
%   Item 2 SEEDS THE CRUX (notes problem 8): the bigger count is not the better rate when the
%          groups are different sizes. Leave both answers on the board.
%   Item 3 seeds Topic 2.1 (VAR-1.D): classify three variables, then ask whether two are related.
\documentclass[12pt]{article}
\usepackage{apstats-article}
\usepackage{apstats-key}   % pulls in -boxes; do NOT also load -boxes

\begin{document}

\pageheader{Unit 2, Lesson 2.1}{Warm-Up --- Answer Key}
% NAMESTRIP: no \namedateperiod here — the name row lives on the cover only.

\vspace{0.28in}

\begin{enumerate}[leftmargin=1.9em, itemsep=0.80in]

  \item \textbf{From Lesson 1.2.} In a survey of 180 students, 45 said they walk to school.
        What is the relative frequency of walkers? Give it as a proportion and a percent.

        \vspace{16pt}
        Proportion: \ans{$45/180 = 0.25$} \qquad Percent: \ans{25\%}

  \item Two sections took the same quiz. In Section A, 12 of 20 students passed. In Section B,
        18 of 40 students passed.

        \vspace{12pt}
        Which section had \emph{more students} pass? \ans{Section B (18)}
        \vspace{16pt}

        Percent who passed: \quad A \ans{60\%} \qquad B \ans{45\%}
        \vspace{16pt}

        Which section did \emph{better} on the quiz? \ans{Section A}

  \item A school survey records three things about each student:
        \vspace{10pt}

        \renewcommand{\arraystretch}{1.5}
        \begin{tabularx}{\linewidth}{@{} Y Y Y @{}}
        \textbf{Phone kept in bedroom (yes/no)} & \textbf{Hours of sleep last night} &
        \textbf{Plays a school sport (yes/no)} \\
        \ans{categorical} & \ans{quantitative} & \ans{categorical} \\
        \end{tabularx}

        \vspace{6pt}
        Label each \textbf{categorical} or \textbf{quantitative}. Then write one question that
        asks whether two of these variables are \emph{related}.\\[18pt]
        \ansline{Do students who keep a phone in the bedroom sleep less than students who}\\[18pt]\ansline{do not? Is playing a sport related to keeping a phone in the bedroom?}

\end{enumerate}

\end{document}
